\section{Introduction}

Password verification is a stream of high-cost point queries: each submitted account--password tuple asks whether an expensive backend hash should run. Modern deployments deliberately use costly PBKDF2, scrypt, or Argon2id parameters~\cite{nist_sp800_63b_2017}; this protects stored passwords, but makes invalid online attempts nearly as expensive as legitimate logins.

\paragraph{Admission-index design problem}
Let \(K_t\) be represented credential keys at epoch \(t\). An admission index returns \textsc{reject} or \textsc{forward}; it is one-sided if every \(q\in K_t\) is always forwarded. Given memory, update limits, risk regions, and a shifting invalid stream, the physical-design objective is to minimize invalid forwards to the backend subject to one-sided correctness and per-region residual-forward budgets. Exact keyed fingerprints and AMQs such as Bloom filters~\cite{broder2004network} can serve as lightweight pre-verifiers, but a single global budget is risk-agnostic: known-password and chosen-password traffic can steer load toward regions where deterministic false positives amplify backend work.

TRAPS is a first-party learned admission layer that keeps learning off the authentication-correctness path. A stable risk oracle routes attempts into regions; a region-local substrate decides only whether to forward to the backend. Learning affects resource allocation rather than authentication correctness: it changes which residual channel and capacity budget a query sees. The oracle never accepts a credential and never rejects a represented credential by itself. Correctness follows from stable features, atomic insertion, and versioned epoch migration; oracle errors only change efficiency.

The abstraction is substrate-agnostic. Exact keyed tags are the strongest fast-verifier point when account-local tag storage is acceptable, while AMQ-backed TRAPS addresses the approximate-membership regime where residual false forwards must be allocated across regions. TRAPS adds capped region allocation, which prevents any routed region from becoming an unbounded steering target, and an exact replay cache that suppresses backend-confirmed false positives. Our Argon2 prototype evaluates controlled credential floods, shift/fallback behavior, LANL status-outcome replay, invalid-heavy trace-conditioned replay, update cost, queueing, and ablations.

\paragraph{Design scope}
The paper evaluates two complementary regimes. Exact keyed tags provide the measured fast-verifier frontier when account-local tag storage is acceptable. AMQ-backed TRAPS studies how to distribute an unavoidable residual false-forward channel across risk regions, replay discoveries, and rebuild epochs. Distribution shift is handled through next-epoch capacity updates, dual-query migration, and drift fallback.

\paragraph{Contributions}
\begin{itemize}
    \item We formulate password-verification floods as a risk-aware one-sided admission-index physical-design problem with memory, update, and regional residual-forward budgets.
    \item We give a capped water-filling allocator with a feasibility certificate and closed-form KKT solution for regional residual channels.
    \item We present AMQ-backed TRAPS: stable learned routing, exact replay suppression, and versioned migration that preserve one-sidedness.
    \item We define a substrate interface for AMQs and keyed fingerprints, including exact-tag rows as the fast-verifier frontier.
    \item We evaluate credential-flood, drift/fallback, LANL-conditioned replay, update, queueing, and ablation workloads.
\end{itemize}
